%%
%% frontmatter/about.tex
%%
%% About this guide: conventions, notation, how to read it.
%%

\chapter*{درباره‌ی این راهنما}
\addcontentsline{toc}{chapter}{درباره‌ی این راهنما}
\label{ch:about}

\section*{ساختار این راهنما}

این راهنما به شش بخش اصلی و پنج پیوست تقسیم شده
است. هر بخش، موضوعی مستقل را پوشش می‌دهد و می‌تواند
به‌عنوان یک مرجع جداگانه استفاده شود.

\subsection*{بخش‌های اصلی}

\begin{description}
    \item[بخش ۱ — شروع به کار:] معرفی \lr{MatinBook}،
    نصب و راه‌اندازی، ایجاد اولین کتاب، و ساختار
    کلی یک کتاب.

    \item[بخش ۲ — اجزای کتاب:] تمام اجزای یک کتاب،
    از جلد و فصل تا منابع و نمایه.

    \item[بخش ۳ — شخصی‌سازی:] آپشن‌ها، رنگ‌ها،
    فونت‌ها، صفحه‌آرایی، تم‌ها و زبان.

    \item[بخش ۴ — مباحث پیشرفته:] ماژول‌ها، تم سفارشی،
    دستورات سفارشی و پروژه‌های بزرگ.

    \item[بخش ۵ — مثال‌های کامل:] سه فصل کامل از
    کتاب‌های ریاضی، برنامه‌نویسی و فیزیک.

    \item[بخش ۶ — مرجع دستورات:] مرجع کامل کلاس،
    محیط‌ها و دستورات.
\end{description}

\subsection*{پیوست‌ها}

\begin{description}
    \item[پیوست الف — عیب‌یابی:] خطاهای رایج و
    راه‌حل آن‌ها.

    \item[پیوست ب — ابزارهای موردنیاز:] نصب و
    پیکربندی \lr{XeLaTeX}، \lr{biber}، \lr{xindy} و
    \lr{Pygments}.

    \item[پیوست پ — منابع:] کتاب‌ها، وب‌سایت‌ها و
    انجمن‌های مفید.

    \item[پیوست ت — تاریخچه:] نسخه‌های \lr{MatinBook}
    و تغییرات آن‌ها.

    \item[پیوست ث — مجوز و مشارکت:] جزئیات مجوز
    و راهنمای مشارکت.
\end{description}

\section*{قراردادهای نگارشی}

در این راهنما، از قراردادهای زیر برای نشان دادن
انواع مختلف محتوا استفاده می‌شود:

\subsection*{نام دستورات}

نام دستورات \lr{LaTeX} با فونت \lr{monospace} و
علامت \lr{\textbackslash} نمایش داده می‌شوند:

\begin{itemize}
    \item \cmd{makecover} — دستور جلد
    \item \cmd{chapter} — دستور فصل
    \item \cmd{cref} — دستور ارجاع
\end{itemize}

\subsection*{نام محیط‌ها}

نام محیط‌ها با فونت \lr{monospace} و بدون
علامت \lr{\textbackslash} نمایش داده می‌شوند:

\begin{itemize}
    \item \env{theorem} — محیط قضیه
    \item \env{minted} — محیط کد
    \item \env{algorithm} — محیط الگوریتم
\end{itemize}

\subsection*{نام بسته‌ها}

نام بسته‌ها با فونت \lr{monospace} نمایش داده
می‌شوند:

\begin{itemize}
    \item \pkg{tcolorbox} — بسته‌ی جعبه‌ها
    \item \pkg{minted} — بسته‌ی کد
    \item \pkg{tabularray} — بسته‌ی جدول
\end{itemize}

\subsection*{نام فایل‌ها}

نام فایل‌ها با فونت \lr{monospace} نمایش داده
می‌شوند:

\begin{itemize}
    \item \file{main.tex} — فایل اصلی
    \item \file{references.bib} — کتاب‌نامه
    \item \file{matinbook.cls} — کلاس اصلی
\end{itemize}

\subsection*{گزینه‌های کلاس}

گزینه‌های کلاس با فونت \lr{monospace} نمایش داده
می‌شوند:

\begin{itemize}
    \item \opt{draft} — حالت پیش‌نویس
    \item \opt{final} — حالت نهایی
    \item \opt{vazirmatn} — فونت وزیرمتن
\end{itemize}

\subsection*{کدها}

کدهای \lr{LaTeX} در بلوک‌های زیر نمایش داده
می‌شوند:

\begin{latin}
\begin{minted}{latex}
\documentclass{matinbook}
\begin{document}
\chapter{|\pc{فصل نمونه}|}
\end{document}
\end{minted}
\end{latin}

\subsection*{خروجی}

خروجی کدها، در صورت نیاز، در بلوک‌های جداگانه
نمایش داده می‌شود. مثال:

\begin{quote}
    \itshape
    این یک متن نمونه است که خروجی کد بالا را
    نشان می‌دهد.
\end{quote}

\section*{نحوه‌ی خواندن این راهنما}

این راهنما به سه روش قابل استفاده است:

\subsection*{روش ۱ — خواندن متوالی}

اگر تازه با \lr{MatinBook} آشنا می‌شوید، توصیه
می‌شود فصل‌ها را به ترتیب بخوانید. هر فصل، بر
پایه‌ی فصل‌های قبلی بنا شده است.

\subsection*{روش ۲ — مرجع سریع}

اگر با \lr{MatinBook} آشنا هستید و فقط می‌خواهید
یک قابلیت خاص را یاد بگیرید، می‌توانید مستقیماً
به فصل مربوطه مراجعه کنید. هر فصل، مستقل قابل
خواندن است.

\subsection*{روش ۳ — جستجو}

برای پیدا کردن یک دستور یا مفهوم خاص، از نمایه‌ی
پایان کتاب استفاده کنید. نمایه، به‌صورت الفبایی
مرتب شده و شامل تمام دستورات، محیط‌ها و مفاهیم
کلیدی است.

\section*{تمرین‌ها}

در پایان هر فصل، مجموعه‌ای از تمرین‌ها ارائه
شده است. توصیه می‌شود این تمرین‌ها را انجام
دهید، زیرا یادگیری \lr{LaTeX} بدون تمرین
عملی امکان‌پذیر نیست.

\section*{نمادها}

در این راهنما، از نمادهای زیر برای تأکید بر
نکات مهم استفاده می‌شود:

\begin{itemize}
    \item \textbf{نکته:} اطلاعات تکمیلی که
    ممکن است مفید باشد.
    \item \textbf{هشدار:} نکاتی که باید به
    آن‌ها توجه کنید تا از خطا جلوگیری شود.
    \item \textbf{مثال:} نمونه‌های عملی از
    مفاهیم.
\end{itemize}

\vspace{1cm}

\begin{flushleft}
    \textbf{متین محمودی} \\
    مهر ۱۴۰۵
\end{flushleft}

\clearpage